\subsection{Ethical considerations}
Business records are less sensitive than personal data, but ER systems in practice also link sole proprietors,
whose business name is often a person's name and whose address is a home. A false merge can attach one trader's
complaints or credit history to another, and the harm falls on the smaller party, who may not know the link
exists. For that reason we favour a precision-weighted metric, keep the model's evidence inspectable through the
logistic regression coefficients and the dashboard, and recommend human review for low-confidence links before they
change any account. The model performs slightly worse on Indian records, which carry more script and spelling
variation; a deployed system should monitor error rates per country and script so that one group of businesses
does not bear most of the mistakes. The challenge data is used only for this academic study; apart from a small demo
sample for the dashboard, it is not redistributed.

\subsection{Limitations}
The study sample keeps the ratio of matched to unmatched records but has 20 times fewer records per country, so
each entity faces fewer look-alike competitors. Blocking recall is therefore higher (99.4\%) than on the full data
(97.8\%), and absolute scores are optimistic compared with the full challenge setting. The comparison between
models is still fair, because all of them see the same candidates and features. The test countries of the
challenge include France, which is absent from training; the lockbox only contains US and Indian records, so our
results say nothing direct about transfer to a new country. Time series methods from the syllabus do not apply,
since the records carry no timestamps.
